\documentclass[11pt,a4paper]{article}
\usepackage[margin=1in]{geometry}
\usepackage[T1]{fontenc}
\usepackage{lmodern}
\usepackage{microtype}
\usepackage{amsmath,amssymb}
\usepackage{xcolor}
\usepackage{hyperref}
\hypersetup{colorlinks=true,linkcolor=blue!50!black,urlcolor=blue!50!black}
\title{Compiling a Parametric C Raytracer to AArch64}
\author{Staged semantic specialization and a fixed-point sphere kernel}
\date{16 August 2026}
\begin{document}
\maketitle
\begin{abstract}
This note expands the semantic-description compiler into a small raytracer compiler instance. A high-level description specifies image dimensions, sphere geometry, and target ABI policy. The compiler specializes that description into AArch64 containing nested render loops, a per-pixel sphere-intersection kernel, an explicit output hook, and a Linux syscall helper. Quoting and unquoting the description produces byte-identical output to direct compilation, demonstrating the self-applicable staging property.
\end{abstract}
\section{Raytracer semantic description}
The source-level description is
\[
D_{ray}=\{
W,H,c_x,c_y,r,\mathsf{entry},\mathsf{syscallABI}\}.
\]
The sample uses $W=64$, $H=48$, center $(32,24)$, and radius $20$. The ray model is orthographic and fixed point: a pixel $(x,y)$ hits when
\[
(x-c_x)^2+(y-c_y)^2<r^2.
\]
The kernel returns intensity $1$ for a hit and $0$ otherwise.

\section{Generated target structure}
The compiler emits a function with the following ABI:
\[
\mathsf{raytrace\_pixel}(x_0,x_1)\to x_0,
\]
where $x_0=x$, $x_1=y$, followed by
\[
\mathsf{pixel\_sink}(x_0,x_1,x_2),
\]
where the arguments are intensity, x, and y. The sink is intentionally explicit: a C variant may implement it using PPM serialization, a framebuffer, or `write(2)` without changing the ray kernel.

The generated control flow is a conventional nested loop:
\begin{verbatim}
.Lrow:
  cmp x19, #48
  b.ge .Lfinish
.Lcolumn:
  cmp x20, #64
  b.ge .Lnextrow
  bl raytrace_pixel
  bl pixel_sink
  b .Lcolumn
\end{verbatim}

The pixel kernel is:
\begin{verbatim}
sub x2, x0, #32
mul x2, x2, x2
sub x3, x1, #24
mul x3, x3, x3
add x2, x2, x3
mov x3, #400
cmp x2, x3
cset x0, lt
ret
\end{verbatim}

\section{Staged compiler generation}
The description is typed code:
\[
\mathsf{quote}:D_{ray}\to\mathsf{Code}(D_{ray}),
\qquad
\mathsf{unquote}:\mathsf{Code}(D_{ray})\to D_{ray}.
\]
The staged compiler is
\[
\mathsf{compileStaged}(\widehat D)=
\mathsf{compile}(\mathsf{unquote}(\widehat D)).
\]
The executable compares direct and staged output as strings. Equality establishes that the semantic description is sufficient to reconstruct the same target compiler instance at the next stage.

\section{Verification}
The implementation is \texttt{raytracer\_arm64.cpp}. It was compiled with:
\begin{verbatim}
g++ -std=c++23 -Wall -Wextra -pedantic -O2 raytracer_arm64.cpp -o raytracer_arm64
./raytracer_arm64 > raytracer_arm64.s
\end{verbatim}
The executable reports:
\begin{verbatim}
raytracer staged equivalence: PASS
\end{verbatim}
Structural checks confirm the presence of `raytrace`, the row and column loops, `raytrace\_pixel`, `pixel\_sink`, the radius-square constant 400, and `sys\_write`.

\section{Limitations}
This is a compiler kernel, not a complete image-producing runtime. It does not yet lower floating point, triangles, recursion, acceleration structures, or a concrete PPM sink. The fixed-point sphere kernel is intentionally small and auditable. The extension path is to add typed vector and image-buffer descriptions, explicit memory effects, and a verified lowering for the output sink while preserving the staged description boundary.
\end{document}
